\chapter{Kidney Biopsy}

\section*{Overview}
A kidney biopsy is a procedure that obtains a small sample of kidney tissue, usually through a needle, to diagnose and guide treatment of medical kidney disease.

\section*{Alternative Names}
Renal biopsy; percutaneous kidney biopsy; renal biopsy specimen

\section*{Principle}
A biopsy needle is advanced into the lower pole of the kidney, typically under ultrasound or CT guidance. Obtained cores are processed for light microscopy, immunofluorescence, and electron microscopy to characterize glomerular, tubular, and interstitial disease.

\section*{History}
Needle biopsy of the kidney was introduced by Iversen and Brun in 1951 and refined with real-time ultrasound guidance in the 1980s, becoming the cornerstone of nephrology diagnosis.

\section*{Why It Is Done}
Performed for nephrotic syndrome, unexplained glomerulonephritis, unexplained acute kidney injury, suspected hereditary kidney disease, and to guide therapy and assess disease activity.

\section*{Is This a Primary or Secondary Test or Procedure}
Primary diagnostic procedure for medical kidney disease.

\section*{How It Is Performed}
Performed in an outpatient or inpatient setting. The patient lies prone, the lower pole is localized by ultrasound, the site is anesthetized, and a spring-loaded biopsy needle obtains two to three cores under real-time guidance. Pressure is held after removal.

\section*{Preparation}
Coagulation studies, blood pressure control, and discontinuation of anticoagulants and aspirin. NPO is usually not required. Blood type and screen are obtained in case of bleeding.

\section*{Contraindications and Precautions}
Uncontrolled hypertension, severe coagulopathy, a solitary kidney, or a bleeding disorder are relative contraindications. Precaution: a small, atrophic or cystic kidney may be unsuitable for biopsy.

\section*{Risks}
Bleeding with macroscopic hematuria, perinephric hematoma, arteriovenous fistula, infection, and rarely the need for embolization or nephrectomy.

\section*{Aftercare and Recovery}
The patient is observed for 6-24 hours with bed rest and serial urinalysis and hematocrit monitoring before discharge.

\section*{Outcome and Follow-up}
Histology classifies glomerular disease (e.g., minimal change, membranous nephropathy, IgA nephropathy, lupus nephritis) and guides immunosuppression.

\section*{Expected Results}
Normal glomeruli, tubules, interstitium, and vessels with no immune deposits.

\section*{Follow-up}
Repeat biopsies may assess treatment response or disease progression; the plan follows the histologic diagnosis.

\section*{References}
\begin{enumerate}
  \item MedlinePlus. Kidney biopsy. U.S. National Library of Medicine; 2025.
  \item Current Medical Diagnosis and Treatment. McGraw-Hill; 2025.
\end{enumerate}

\clearpage
